\section{WC 21/7/25}
\subsection*{Overview:}
This week I planned out my Product demonstration video that will serve as the video artifact i submit along with my EP Report (As the website alone cannot be submitted)

\subsection*{Achievements, Thoughts and Ideas:}
I started the week by conducting some research into how other competing platforms present their offerings in a variety of product videos: \nameref{research:ProductShowcaseVideos}.

I found that the most effective videos were those that had a clear structure and narrative whilst remaining very calm and focused on success, showcasing the product's features in a way that was engaging and easy to understand. I also noted the importance of high-quality visuals and sound design in creating a professional look and feel
.
I then began to plan my own video, I wanted to create a video that was engaging and informative, but also had a clear narrative and structure. I decided to use a combination of voiceover, visuals and sound effects to create a cohesive and engaging experience for the viewer. I already have some experience with video editing and production thanks to some other extracurricular events i volunteer within (such as cinema and theatre), so I felt confident in my ability to leverage these skills to create a high-quality video.

Firstly I wanted to plan out my initial ideas regarding the videos structure, here is my work in progress script for the video including the various scenes and voiceovers:

\textbf{Scene 1: Introduction}
\textit{set in a dark messy bedroom, camera faces a desk with a monitor to the right and glass sphere on left}
\begin{itemize}
    \item [0:00] \textbf{Voiceover:} "Do you ever feel like your notes, thoughts and ideas are being forgotten"
    \item \textbf{Visuals:} a sphere of glass with glowing purples and blues is fading / pulsing in brightness
    \item [0:00] \textbf{SFX:} pulsing energy sound effect
    \item [0:05] \textbf{Voiceover:} "that they are lost in the noise of every day life"
    \item \textbf{Visuals:} the camera pans slightly to show the messy room, the light from the sphere dies, a moving beam of light shines on the room (train tracks etc)
    \item [0:10] \textbf{SFX:} "train rumbling, tracks vibrating representing lost trains of thought"
    \item [0:10] \textbf{Voiceover:} "but what if that could change"
    \item \textbf{Visuals:} the sphere lights up again, the camera pans back to the desk
    \item [0:15] \textbf{SFX:} "strong sharp energy 'click' sound"
    \item [0:15] \textbf{Voiceover:} "what if your greatest assets could be reunited again."
    \item \textbf{Visuals:} spheres all emerge from the depths of the room, each one representing a different thought or idea, they are all glowing and pulsing
    \item [0:20] \textbf{SFX:} "rustling sounds and energy pulsing"
    \item [0:20] \textbf{Voiceover:} "Introducing NeuroNote"
    \item \textbf{Visuals:} the spheres all come together to form a glowing NeuroNote logo, the camera zooms in on the logo
    \item [0:20] \textbf{Voiceover:} "A portal into your mind"
    \item \textbf{Visuals:} spheres stretch into a portal like "stars wars hyperspace" effect, the camera zooms into the portal and into the computer monitor
\end{itemize}

\textbf{Scene 2: High level benefits}
\textit{set within the computer monitor, simple 2d animation style}
\begin{itemize}
    \item [0:30] \textbf{Voiceover:} "Whether you are out for a run, lying in bed or at your desk"
    \item \textbf{Visuals:} the 3 people appear on screen, each one representing a different use case for NeuroNote
    \item [0:30] \textbf{Voiceover:} "You will always be connected"
    \item \textbf{Visuals:} the three people are all dragged towards the central NeuroNote logo and begin to orbit slowly turning into the notes graph view
    \item [0:35] \textbf{Voiceover:} "Welcome to the graph view"
    \item [0:40] \textbf{Visual:} orbiting people turn into a graph view, the camera zooms in on the graph view
    \item [0:40] \textbf{Voiceover:} "A visual representation of your thoughts and ideas, easily navigatible, searchable and always at your fingertips. Positioned at the very heart of NeuroNote"
    \item [0:45] \textbf{Visuals:} the camera zooms out to show the entire graph view within the application dashboard, the rest of the page beats on the word heartbeat.
    \item [0:45] \textbf{SFX:} "heartbeat sound effect"
    \item \item [0:50] \textbf{Voiceover:} "the graph view beautifully visualises your thoughts and ideas, allowing you to see the connections between them and how they relate to each other. Notes connect to tags and tags to other tags creating a web of knowledge that is easy to navigate and understand. With powerful search and filtering capabilities, you can quickly find the information you need, when you need it."
    \item 
    \item [0:50] \textbf{Voiceover:} "With NeuroNote, you can say goodbye to the chaos of you archaic note-taking systems and breath new life into your thoughts with the press of a button"
    \item [0:55] \textbf{Visuals:} the software makes a new document type note adnd the camera ripples in
\end{itemize}

\textbf{Scene 3: Features}
\textit{set within the computer monitor, simple 2d animation style}
\begin{itemize}
    \item [1:00] \textbf{Voiceover:} "Embrace the editor, a powerful tool for creating and editing your notes. packed full of features to help you craft your perfect memory. "
    \item [1:05] \textbf{Visuals:} the camera zooms in on the note creation screen, the editor appears on screen with a variety of features highlighted
    \item [1:10] \textbf{Voiceover:} "By adding a huge range of multimedia 'blocks' to your notes you can create rich, engaging and informative notes that are easy to read and understand"
    \item [1:15] \textbf{Visuals:} the camera zooms in on the note creation screen, the user adds a variety of blocks to the note, each one representing a different type of content (text, image, video, audio, etc.)
    \item [1:20] \textbf{Voiceover:} "With NeuroNote's notes, you can record ideas that are truly unique and tailored to your needs"
    \item [1:20] \textbf{Voiceover:} "but what if you find yourself craving the freedom of pen and paper? Enter the Whiteboard - your handwriting-fueled sidekick. Draw diagrams, jot spontaneous thoughts, or map ideas just like you would on a real paper — all while staying connected to your digital graph." 
    \item 
\end{itemize}

textbf{Scene 4: call to action}
\textit{set within the computer monitor, simple 2d animation style}
\begin{itemize}
    \item [1:30] \textbf{Voiceover:} "So what are you waiting for? Join the NeuroNote revolution today and start taking control of your thoughts and ideas"
    \item [1:35] \textbf{Visuals:} the camera zooms out to show the entire application, the NeuroNote logo appears on screen
    \item [1:40] \textbf{Voiceover:} "Visit our website at neuronote.space to get started for free!"
    \item [1:45] \textbf{SFX:} "upbeat music fades out"
    \item [1:50] \textbf{Voiceover:} "NeuroNote, the future of note-taking"
\end{itemize}




